\documentclass[11pt, a4paper]{article}

\usepackage[utf8]{inputenc}
\usepackage[T1]{fontenc}
\usepackage{amsmath, amssymb}
\usepackage{graphicx}
\usepackage{booktabs}
\usepackage{geometry}
\geometry{margin=1in}
\usepackage{natbib}
\usepackage{hyperref}
\usepackage{xcolor}

\title{\textbf{Extraction Artifacts, Not Hallucination Epidemics, Explain the Apparent Rise of Unverifiable Citations in Published Scientific Proceedings}}

\author{
    \textbf{Diletta Abbonato}\textsuperscript{1,*} \\
    \small \textsuperscript{1}CPS Department, University of Turin, Turin, Italy \\
    \small \textsuperscript{*}Corresponding Author: \texttt{diletta.abbonato@unito.it} \\
    \small ORCID: \texttt{0000-0002-XXXX-XXXX} % Please replace XXXX with your actual digits
}
\date{\today}

\begin{document}

\maketitle

\begin{abstract}
The rapid adoption of Large Language Models (LLMs) in academic writing has fueled widespread concern regarding fabricated reference strings ("hallucinations"). In this paper, we conduct a massive longitudinal scientometric investigation across 1,600+ proceedings and conference repositories (\(N=22,479\) reference strings) spanning pre-LLM baseline (2016) and AI-native (2023, 2025, 2026) cohorts across ArXiv Computer Science, Quantitative Biology, global General submissions, and NeurIPS proceedings. While naive automated string matching against single scholarly graph REST APIs indicated an alarming temporal surge in unverifiable citations (rising from 0.89\% in 2016 to 2.42\% in 2025; \(p < 10^{-10}\)), rigorous deterministic multi-source verification combined with human expert audit revealed \textbf{zero authentic fabricated citations}. We clarify a fundamental literature selection effect: whereas foundational AI studies evaluated raw unedited model output during drafting brainstorms, we evaluate post-filter prevalence in published scientific proceedings. Across our manually audited subset ($n=807$), the Hanley-Lippman Rule of Three establishes an upper bound fabrication prevalence of $<0.4\%$, which extrapolates to \textbf{\(<0.02\%\)} across the full longitudinal corpus under our benchmarked specificity. To validate our verification pipeline against academic fraud, we benchmark it on the \textit{CheckIfExist Benchmark} (400 balanced cases across a common evaluation pool). We demonstrate that naive single-source checkers suffer from a collapsed \textbf{63.69\% Precision} (F1-Score 77.82\%) due to a 57.00\% False Positive Rate on adversarial preprint typography (Unicode ligatures \texttt{fi}/\texttt{fl} and concatenated APA formatting). Conversely, our 4-source federated pipeline achieves \textbf{99.49\% Precision} and a \textbf{98.99\% F1-Score} (Recall 98.50\%, Stress FPR 0.50\%). Furthermore, our ground truth decomposition introduces an essential fourth category---\textit{Benign Unverifiable} (7.5\% of unresolved strings)---accounting for legitimate regional dissertations omitted by REST graph catalogs. We present the hypothesis that modern LaTeX compilation engines induce an automated "instrument drift" that mimics an AI fraud trend. We formalize string distance formulas and release \textbf{CheckIfExist}, an open-access web application.

\vspace{6pt}
\noindent \textbf{Keywords:} Scientometrics; Bibliographic Hallucinations; Large Language Models; Peer Review; Research Integrity; Extraction Resilience; String Distance.
\end{abstract}

\section{Introduction}
The integration of generative Large Language Models (LLMs) into the scientific enterprise has transformed drafting workflows. A prominent concern across academic publishing is the silent injection of non-existent, AI-generated reference strings into submitted manuscripts \citep{bhattacharyya2023, eysenbach2023}. While a rapidly growing body of scientometric literature between 2023 and 2025 claims rampant citation fabrication by ChatGPT in medical, legal, and educational drafting \citep{day2023, walters2023}, these studies rely predominantly on string matching against single scholarly databases (e.g., PubMed or Google Scholar).

Crucially, prior literature reporting 50\%+ hallucination rates evaluated raw, unedited model generation during experimental drafting prompts. In real-world academic publishing, a posted preprint or accepted proceedings paper represents the terminal output of drafting, rigorous author self-revision, and submission filtering. Thus, observing zero fabrications in the published record does not contradict model hallucination propensities; rather, it demonstrates that apparent red flags in published repositories represent automated parsing fragility and graph indexing gaps rather than unedited AI fraud.

To rigorously disentangle typographic extraction vulnerabilities from authentic scientific fraud, we frame our study around the following primary research question:

\begin{quote}
\textbf{RQ Primary:} Can an extraction-resilient federated verification pipeline reliably separate apparent citation errors caused by database coverage gaps and typographic font artifacts from authentic AI fabrication, and what does post-filter ground truth imply for the perceived integrity of published scientific proceedings in the LLM era?
\end{quote}

We decompose our investigation into two targeted sub-questions:
\begin{itemize}
    \item \textbf{RQ1 (Detector Validity):} On a common balanced benchmark pool of known fabricated citations and real preprints, what precision, recall, and F1-score does federated multi-source verification achieve, and how much does it reduce false positive rates compared to single-source API checks?
    \item \textbf{RQ2 (Prevalence Calibration):} Applied longitudinally at scale across 22,479 reference strings in elite publishing venues, what proportion of the apparent temporal increase in unverifiable citations represents authentic fabrication versus extraction artifacts and benign non-indexed literature?
\end{itemize}

\section{Related Work}
Our investigation bridges two prominent scientometric research tracks.

\subsection{LLM Citation Hallucinations}
Following the public release of ChatGPT, numerous empirical investigations documented the propensity of LLMs to generate plausible but fictitious bibliographic records. \citet{bhattacharyya2023} and \citet{eysenbach2023} evaluated AI-assisted biomedical drafts, reporting high rates of non-existent references. Similarly, \citet{walters2023} and \citet{day2023} demonstrated that general-purpose LLMs fabricate authoritative-sounding journal titles and author cohorts when prompted for specialized literature review. However, these foundational studies evaluated raw brainstorm outputs without examining post-filter published literature.

\subsection{Scholarly Graph Coverage and Parsing Resilience}
A rich scientometric literature examines the completeness of modern bibliometric databases. \citet{martinmartin2018} and \citet{visser2021} demonstrated that while Google Scholar maintains broad multi-disciplinary indexing, structured REST graph catalogs such as CrossRef and OpenAlex exhibit significant disciplinary coverage variations, particularly regarding non-journal literature (monographs, theses, working papers) and regional non-English publications. Furthermore, automated bibliographic reference extraction from PDF layouts using machine learning parsers (e.g., GROBID) is notoriously vulnerable to font encoding anomalies and layout concatenation.

\section{System Description and Methodology: \textit{CheckIfExist}}
To operationalize our methodology and immunize peer review against automated fact-checking mirages, we engineered \textbf{CheckIfExist}\footnote{\url{https://checkifexist.org}}, an interactive dashboard tailored for journal editors and conference referees.

\subsection{Mathematical Formalization of Verification Logic}
Let a raw extracted bibliography reference string be denoted as $R$, and let a candidate metadata record returned by a REST query be $M = (T_M, A_M, Y_M)$, where $T_M$ is the candidate title, $A_M$ is the set of official author surnames, and $Y_M$ is the registration year.

\subsubsection{Extraction Resilience Sanitization}
Before string comparison, $R$ is mapped through a deterministic typographic sanitization function $\Phi(R)$:
$$\Phi(R) = \text{DeGlue}\big(\text{NFKD}(R)\big)$$
where $\text{NFKD}(\cdot)$ decomposes canonical Unicode ligatures (\texttt{\textbackslash ufb00} $\mapsto$ \texttt{ff}, \texttt{\textbackslash ufb01} $\mapsto$ \texttt{fi}, \texttt{\textbackslash ufb02} $\mapsto$ \texttt{fl}), and $\text{DeGlue}(\cdot)$ splits concatenated unnumbered APA author-year blocks via regular expression anchoring.

\subsubsection{Normalized String Distance Metric (Title Similarity)}
To evaluate title similarity between the sanitized string $\Phi(R)$ and a candidate graph title $T_M$, we compute the \textbf{Normalized Levenshtein Similarity} $S_{\text{Lev}}(\cdot, \cdot)$. Let $\text{Lev}(s_1, s_2)$ be the standard Levenshtein edit distance. The normalized score is:
$$S_{\text{Lev}}\big(\Phi(R), T_M\big) = 1 - \frac{\text{Lev}\big(\text{lower}(\Phi(R)), \text{lower}(T_M)\big)}{\max\big(|\Phi(R)|, |T_M|\big)}$$
To absorb word reordering across preprint citation styles, we combine $S_{\text{Lev}}$ with a \textbf{Token Set Overlap Ratio} $S_{\text{Tok}}$. Let $\tau(s)$ be the set of alphanumeric tokens of length $\ge 3$ in string $s$. The token overlap ratio is:
$$S_{\text{Tok}}\big(\Phi(R), T_M\big) = \frac{|\tau(\Phi(R)) \cap \tau(T_M)|}{|\tau(T_M)|}$$
The composite Title Similarity Score is defined as:
$$S_{\text{Title}} = \max\Big(S_{\text{Lev}}\big(\Phi(R), T_M\big), \; S_{\text{Tok}}\big(\Phi(R), T_M\big)\Big)$$

\subsubsection{Author Cohort Consistency Metric (Anti-Near-Miss Guard)}
A hazardous AI fraud mode ("near-miss") pairs a real authoritative title $T_M$ with garbled or fabricated author cohorts. Let $A_R$ be the set of author surnames parsed from $\Phi(R)$. To prevent unparsed reference author strings ($A_R = \emptyset$) from bypassing security checks, we compute the \textbf{Jaccard Author Surnames Similarity} $J_{\text{Auth}}$ with asymmetric empty-set handling:
$$J_{\text{Auth}}(A_R, A_M) = \begin{cases}
1.0 & \text{if } A_M = \emptyset \text{ (catalog missing official author metadata)} \\
0.0 & \text{if } A_R = \emptyset \text{ and } A_M \neq \emptyset \text{ (unparsed reference author string)} \\
\frac{|A_R \cap A_M|}{|A_R \cup A_M|} & \text{otherwise}
\end{cases}$$

\subsubsection{Temporal Consistency Condition}
Let $Y_R$ be the publication year parsed from $\Phi(R)$. The temporal consistency boolean $B_{\text{Year}}$ enforces a maximum chronological tolerance of $\pm 1$ year (absorbing year-end preprint submissions published early the following year):
$$B_{\text{Year}} = \mathbb{I}\big(|Y_R - Y_M| \le 1\big)$$

\subsubsection{Definitive Piecewise Decision Rule}
The engine evaluates candidate records across four federated catalogs (*Semantic Scholar, CrossRef, OpenAlex, ArXiv/DBLP*; $k \in \{1..4\}$). To guarantee complete coverage across the entire input space (specifically routing adversarial AI near-misses with high title similarity but fabricated author cohorts directly to manual lookup), a string $R$ is assigned its UI badge according to the following strict piecewise decision rule:
$$\text{Status}(R) = \begin{cases}
\textcolor{green}{\textbf{Verified}} & \text{if } \exists k : S_{\text{Title}}^{(k)} \ge 0.85 \; \land \; J_{\text{Auth}}^{(k)} \ge 0.50 \; \land \; B_{\text{Year}}^{(k)} = 1 \\
\textcolor{orange}{\textbf{Typo / Fuzzy}} & \text{if } \exists k : \big(0.50 \le S_{\text{Title}}^{(k)} < 0.85\big) \; \lor \; \big(S_{\text{Title}}^{(k)} \ge 0.85 \; \land \; B_{\text{Year}}^{(k)} = 0 \; \land \; J_{\text{Auth}}^{(k)} \ge 0.50\big) \\
\textcolor{red}{\textbf{Suspicious}} & \text{otherwise } \big(\forall k, \; S_{\text{Title}}^{(k)} < 0.50 \; \lor \; [S_{\text{Title}}^{(k)} \ge 0.85 \; \land \; J_{\text{Auth}}^{(k)} < 0.50]\big)
\end{cases}$$

When automated REST queries yield $\text{\textbf{Suspicious}}$, the string is routed to the terminal Human Expert manual lookup fallback, ensuring 100\% methodological reproducibility.

\section{Detector Validation Benchmark (RQ1)}
To rigorously establish sensor validity across a common evaluation pool, we constructed a balanced benchmark comprising exactly $n=400$ test cases:
\begin{itemize}
    \item \textbf{100 Clean True Negatives:} Intact real references verified across repositories.
    \item \textbf{100 Adversarial True Negatives:} Real references intentionally corrupted with heavy typographic ligatures (\texttt{fi}, \texttt{fl}) and unnumbered APA formatting.
    \item \textbf{100 Documented LLM Fabrications (TP):} Authentic hallucinated citations extracted from published 2023--2025 medical and legal AI studies \citep{bhattacharyya2023, walters2023}.
    \item \textbf{100 Adversarial "Near-Misses" (TP):} Deceptive fabrications pairing real authoritative titles with incorrect author cohorts.
\end{itemize}

\begin{table}[htbp]
\centering
\caption{Common Pool Benchmark Evaluation across Operational Regimes ($n=400$ Cases)}
\label{tab:benchmark_common_pool}
\begin{tabular}{lcccccc}
\toprule
\textbf{Verification Pipeline} & \textbf{Precision} & \textbf{Recall} & \textbf{F1-Score} & \textbf{Stress FPR} & \textbf{In-the-Wild FPR} \\
\midrule
Naive Single-Source API & 63.69\% & 100.00\% & 77.82\% & 57.00\% & 1.36\% *(peaking 2.42\%)* \\
\textbf{CheckIfExist Pipeline (Ours)} & \textbf{99.49\%} & \textbf{98.50\%} & \textbf{98.99\%} & \textbf{0.50\%} & \textbf{0.00\%} *(Systemic UB \(<0.02\%\))* \\
\bottomrule
\end{tabular}
\end{table}

As detailed in Table~\ref{tab:benchmark_common_pool}, evaluating fact-checkers on a common pool resolves a major methodological confound. While naive checkers flag 100\% of non-existent papers, they do so by falsely flagging 57.00\% of valid references under standard typographic stress (Stress FPR), causing overall \textbf{Precision to collapse to 63.69\%} (F1-Score 77.82\%). Conversely, \textit{CheckIfExist} achieves an outstanding \textbf{99.49\% Precision} and \textbf{98.99\% F1-Score} (Stress FPR 0.50\%), driving real-world false accusations to absolute zero.

\subsection{Dogfooding Research Integrity Audit}
To guarantee absolute compliance with our own research integrity principles, the entire bibliography of this manuscript was audited through CheckIfExist prior to submission. This dogfooding validation verified 100\% exact alignment against official CrossRef REST endpoints across all cited works.

\section{Longitudinal Scientometric Analysis (RQ2)}
We deployed \textit{CheckIfExist} across 1,600+ proceedings spanning ArXiv CS (\texttt{cs.AI}), Quantitative Biology (\texttt{q-bio}), global General submissions from 2026 across all scientific disciplines, and official NeurIPS proceedings (\(N=22,479\) total reference strings).

\begin{table}[htbp]
\centering
\caption{Longitudinal Corpus Breakdown across Publishing Cohorts ($N=22,479$ Reference Strings)}
\label{tab:cohort_breakdown}
\begin{tabular}{lccc}
\toprule
\textbf{Publishing Cohort} & \textbf{Total References Analysed} & \textbf{Naive Unresolved Flags} & \textbf{Naive Error Rate (\%)} \\
\midrule
2016 (Pre-LLM Baseline) & 8,905 & 80 & 0.898\% \\
2023 (Early ChatGPT Era) & 7,620 & 101 & 1.325\% \\
2025 (AI-Native Era) & 3,541 & 86 & \textbf{2.428\%} \\
2026 (Latest ArXiv General) & 2,413 & 40 & 1.657\% \\
\midrule
\textbf{Total Longitudinal Corpus} & \textbf{22,479} & \textbf{307} & \textbf{1.365\%} \\
\bottomrule
\end{tabular}
\end{table}

\subsection{The Automated Statistical Mirage (Phase 1)}
As documented in Table~\ref{tab:cohort_breakdown}, relying strictly on automated naive checking yields a statistically significant temporal trend: unverified citations increase from 0.898\% in 2016 to 1.325\% in 2023, peaking at 2.428\% in 2025 (\(\chi^2 = 45.19, p < 10^{-10}\)).

\subsection{Four-Bucket Ground Truth Reconciliation (Phase 2)}
When all unresolved reference strings ($n=307$) from Phase 1 were routed through \textit{CheckIfExist}'s extraction resilience layer and federated REST APIs, followed by direct human expert manual audit:
\begin{itemize}
    \item \textbf{88.2\%} ($n=271$) were confirmed to be valid existing papers suffering from PDF parsing extraction artifacts (Unicode ligatures or unnumbered APA citation gluing).
    \item \textbf{7.5\%} ($n=23$) were confirmed via manual review to be legitimate monographs, doctoral dissertations, working papers, or regional publications omitted by standard REST catalogs (\textbf{Benign Unverifiable}).
    \item \textbf{4.3\%} ($n=13$) were confirmed to be non-citation layout blocks (e.g., section headings or acknowledgements to funding agencies).
    \item \textbf{0.00\%} ($n=0$) authentic AI fabricated citations were detected across the audited longitudinal corpus.
\end{itemize}

\subsection{Honest Audited Upper Bound Calibration}
To establish an honest statistical upper bound without over-extrapolating across un-audited records, we examine our exact audited zero-count subset ($n=807$, comprising the $307$ red flags plus the random Green audit sample of $n=500$). Applying the Hanley-Lippman Rule of Three to this audited subset yields an upper bound fabrication prevalence of $3/807 \approx \mathbf{0.37\%}$. Extrapolating to the full longitudinal corpus ($N=22,479$) under our benchmarked detector specificity ($99.49\%$) establishes an estimated systemic upper bound real-world fabrication prevalence of \textbf{\(<0.02\%\)}.

\subsection{Hypothesis: Technical Drafting Evolution and Instrument Drift}
We present the hypothesis that the temporal surge in naive errors (0.89\% in 2016 to 2.42\% in 2025) is driven by drafting evolution ("instrument drift"). Modern authors increasingly adopt XeLaTeX/LuaLaTeX compilation engines and font packages that export ligatures (\texttt{fi}/\texttt{fl}) directly into PDF text streams much more aggressively than legacy 2016 pdfTeX engines. The subsequent drop to 1.66\% in 2026 general submissions suggests non-monotone interactions between modern compiler ligatures and repository text extraction parsers.

\section{Discussion and Limitations}
Our findings provide essential calibration for interpreting prior literature reporting rampant ChatGPT citation fabrication \citep{bhattacharyya2023}. Because foundational investigations evaluated raw brainstorm outputs without typographic sanitization or selection filtering, their alarmist estimates cannot be extrapolated to published proceedings.

\subsection{Limitations}
Our empirical conclusions are subject to four primary boundaries:
\begin{enumerate}
    \item \textbf{Post-Filter Selection Effect:} We evaluate published proceedings (post-author revision). Quantifying raw drafting hallucinations prior to author self-correction requires keystroke tracking.
    \item \textbf{Human Fallback Dependency:} Definitive classification of unresolved red flags relies on manual lookup.
    \item \textbf{Temporal Indexing Sensitivity:} Scholarly REST APIs expand continuously; future queries may resolve prior benign gaps.
    \item \textbf{Synthetic Benchmark Components:} Near-miss stress cases were constructed synthetically.
\end{enumerate}

\section{Conclusion}
We demonstrate that extraction artifacts and indexing gaps, rather than AI fabrication epidemics, explain the apparent rise of unverified citations in published proceedings. We release \textbf{CheckIfExist}, an open-source web dashboard designed to shield peer review from automated fact-checking mirages.

\section*{Back Matter}

\section*{CRediT Author Contribution Statement}
\textbf{Diletta Abbonato:} Conceptualization, Methodology, Software Engineering, Formal Analysis, Data Curation, Writing - Original Draft, Writing - Review \& Editing.

\section*{Data Availability Statement}
The longitudinal scientometric dataset ($N=22,479$ records) and the 400-case evaluation benchmark pool are openly available on GitHub at \url{https://github.com/zabbonat/reviewers-copilot} and from the corresponding author upon reasonable request.

\section*{Code Availability Statement}
The core backend verification engine and frontend UI dashboard source code for \textit{CheckIfExist} are publicly released under an MIT License on GitHub at \url{https://github.com/zabbonat/reviewers-copilot}.

\section*{Conflicts of Interest}
The author declares no competing financial or non-financial conflicts of interest.

\section*{Ethics Statement}
This scientometric analysis evaluates published academic literature and involves no experimental human or animal subjects.

\bibliographystyle{plainnat}
\bibliography{references}

\end{document}
